Large language models (LLMs) often struggle to use tools reliably in domain-specific settings, where APIs may be idiosyncratic, under-documented, or tailored to private workflows. This highlights the need for effective adaptation to task-specific tools. We propose \shortname, a neuro-symbolic approach for LLM adaptation based on dynamic rule injection. Compact, interpretable rules are distilled from failure traces and injected into the prompt during inference to improve task performance. These rules are proposed by the LLM itself and consolidated using a Minimum Description Length (MDL) objective that favors generality and conciseness. Each rule is stored in both natural language and a structured symbolic form, supporting efficient retrieval at inference time. Experiments on tool-use benchmarks show that this approach improves accuracy on both seen and unseen tools without modifying LLM weights. It outperforms prompting-based adaptation methods and complements finetuning. Moreover, rules learned from one LLM can be reused to improve others, including long reasoning LLMs, highlighting the portability of symbolic knowledge across architectures.%\footnote{The code will be released upon acceptance.}.
\footnote{\emph{Proceedings of the 64th Annual Meeting of the Association for Computational Linguistics} (Volume 1: Long Papers), pages 34631–34646
July 2-7, 2026 (ACL 2026). Copyright 2026
by the authors.}

